\section{Polynomial pairings on a punctured affine space}
\label{sec:pairing}

The pairing comes from low-degree polynomials that distinguish the points
to be deleted. Their pairwise products sum to zero on every full affine
line and on the whole affine space. At the deleted points, however, the
evaluation vectors are coordinate vectors, so their contribution is the
identity matrix. Removing these points leaves its negative as the total
pairing on the remaining space.

The interpolation argument does not depend on how the retained full lines
are arranged. Regularity and girth enter as separate geometric hypotheses
in the proposition below. Section~\ref{sec:hypergraph} will supply them,
after Sections~\ref{sec:criterion} and~\ref{sec:packing} explain their
algebraic use.

We use the usual lower-set Newton interpolation argument; compare
Dyn and Floater~\cite{DynFloater2014}. The triangular
evaluation proof below works directly over the finite field needed here.
The later products of interpolants and their restrictions to full affine
lines are the features related to the coding constructions of
Golowich--Lin \cite[Sections~2.1.2--2.1.3]{GolowichLin2024} and
Dinur--Liu--Zhang \cite[Sections~3.4--3.5]{DinurLiuZhang2025}.
We need no coding theorem from either work: the degree bound and all line
identities used in this paper are proved below.

\begin{lemma}[Interpolation on an index simplex]
\label{lem:interpolation}
Let $Q$ be a field of order $q$, enumerate its elements as $a_0,\ldots,a_{q-1}$, and let $0\le d<q$. For $n\ge1$, set
\[
S_d=\{(a_{k_1},\ldots,a_{k_n}): k_i\ge0,\ \textstyle\sum_i k_i\le d\}\subset Q^n.
\]
Evaluation on $S_d$ is an isomorphism from the $Q$-vector space of polynomials of total degree at most $d$ to $Q^{S_d}$.
\end{lemma}
\begin{proof}
For each tuple $k=(k_1,\ldots,k_n)$ of total sum at most $d$, consider
\[
N_k(X_1,\ldots,X_n)=\prod_{i=1}^n\prod_{0\le u<k_i}(X_i-a_u).
\]
Its evaluation at the point with index tuple $j$ vanishes if $j_i<k_i$ for any $i$. Thus, with tuples in lexicographic order, the square evaluation matrix is triangular. Its diagonal entry indexed by $k$ is
\[
\prod_{i=1}^n\prod_{0\le u<k_i}(a_{k_i}-a_u)\ne0.
\]
Since $d<q$, there is no additional coordinate restriction on these tuples. Their number is $\binom{d+n}{n}$, the dimension of the polynomial space. The displayed polynomials therefore give an invertible evaluation map.
\end{proof}

\begin{lemma}[Line sums]
\label{lem:line-sums}
Let $n\ge1$. If $P$ is a polynomial on $Q^n$ of total degree at most $q-2$, then its sum on every full affine $Q$-line is zero. Its sum on $Q^n$ is also zero.
\end{lemma}
\begin{proof}
On a line, write its restriction as a polynomial in a parameter $z\in Q$ of degree at most $q-2$. The sum of the constant monomial is $q=0$ in $Q$. If $1\le k<q-1$, choose $c\in Q^\times$ with $c^k\ne1$, using cyclicity of $Q^\times$. Multiplication by $c$ permutes $Q$, so
\[
\sum_{z\in Q}z^k=c^k\sum_{z\in Q}z^k
\]
forces the sum to vanish. More explicitly, $(1-c^k)\sum_z z^k=0$ and $1-c^k\ne0$. This proves the line assertion. Partition $Q^n$ into parallel lines to obtain the last assertion.
\end{proof}

\begin{proposition}[The pairing supplied by holes]
\label{prop:pairing}
Let $q\ge4$ be a power of two, let $Q$ be a field of order $q$, fix an enumeration of $Q$, and let $n\ge1$. Put $d_0=(q-2)/2$, $S=S_{d_0}\subset\mathcal A=Q^n$, and $N_{\rm pt}=q^n$. Let $\ell$ be an odd prime and $D=\ell+1$. Suppose that $O\subset S$ and that $\mathcal F$ is a family of full affine $Q$-lines contained in $V=\mathcal A\setminus O$, with
\begin{equation}
\label{eq:hole-size}
\frac{2N_{\rm pt}}{D^2}\le |O|\le\frac{7N_{\rm pt}}{D^2}.
\end{equation}
Suppose also that every point of $V$ lies on $D$ members of $\mathcal F$ and that the incidence graph has girth at least $12$. If $K$ is a finite extension of $Q$ containing a primitive $\ell$th root of unity, then, with $m=|O|$ and $t=|V|$, there are columns $x_v\in K^m$, $v\in V$, such that
\[
0<m<t\le m\ell^2,\qquad
\sum_{v\in V}x_vx_v^{\mathsf T}=-I_m,\qquad
\sum_{v\in f}x_vx_v^{\mathsf T}=0\quad(f\in\mathcal F).
\]
\end{proposition}
\begin{proof}
Interpolation is on the whole set $S$, although only the subset $O$ is
deleted. For each $o\in O$, Lemma~\ref{lem:interpolation} gives a unique
polynomial $f_o$ of total degree at most $d_0$ whose values on $S$ are
the indicator function of $o$. Order $O$ and define, for all
$v\in\mathcal A$,
\[
x_v=(f_o(v))_{o\in O}\in Q^m\subset K^m,
\qquad W_v=x_vx_v^{\mathsf T}.
\]
Each entry of $W_v$ is the value at $v$ of a polynomial of total degree at most $2d_0=q-2$. Lemma~\ref{lem:line-sums} therefore gives
\[
\sum_{v\in f}W_v=0\quad(f\in\mathcal F),
\qquad \sum_{v\in\mathcal A}W_v=0.
\]
For $o\in O$, $x_o$ is its coordinate unit vector. Hence $\sum_{o\in O}W_o=I_m$, and subtraction yields
\begin{equation}
\label{eq:polynomial-pairing}
\sum_{v\in V}W_v=-I_m,
\qquad \sum_{v\in f}W_v=0\quad(f\in\mathcal F).
\end{equation}

It remains to check the sizes. Since $D\ge4$, the upper bound in \eqref{eq:hole-size} is strictly less than $N_{\rm pt}/2$. The lower bound is positive. Thus $0<m<t$, and
\[
\frac{t}{m}=\frac{N_{\rm pt}}{m}-1
\le\frac{D^2}{2}-1\le(D-1)^2=\ell^2,
\]
where the last inequality is equivalent to $(D-2)^2/2\ge0$. These are all the claimed identities and inequalities.
\end{proof}

\begin{remark}
The simultaneous identities in \eqref{eq:polynomial-pairing} are compatible with regularity because $D=\ell+1$ is zero as a scalar in characteristic two. Summing the line identities multiplies the total point sum by $D$, rather than recovering that total point sum.
\end{remark}
